\section{Introduction}

Why do some industries innovate rapidly while others stagnate? When
innovation concentrates geographically, as in semiconductors, software, and
biotechnology, what role do cross-firm knowledge spillovers play?
\citet{igami_estimating_2017} answers a version of the first question for
the hard-disk-drive industry, where the central friction is cannibalization.
Adoption of a new generation forces an incumbent to compete in the new
market against its own old-generation product, while pure entrants face no
such conflict. Incumbents therefore adopt more slowly than entrants, and
they lose share as the new generation takes over: the innovator's dilemma.\\

We extend the Igami model in two directions. First, we introduce regional
heterogeneity: $R = 3$ regions, each with its own roster of old incumbents,
adopters, new entrants, and potential entrants. Second, we add bloc-pooled
agglomeration spillovers, in which the marginal cost of adopting new
technology in region $r$ falls in the total number of adopters in $r$'s
spillover bloc. The default calibration assigns each region to its own
singleton bloc, recovering purely local spillovers and nesting the Igami
baseline \citep{carlino_kerr_2015_agglomeration}. \textcolor{red}{[Pedro]
cite a paper documenting that agglomeration effects are quantitatively
important.}\\

The paper makes three contributions. Methodologically, we show that
\citeauthor{igami_estimating_2017}'s sequential-move solver extends to
multiple regions. Within each period, stages move in the order old $\to$
both $\to$ new $\to$ potential entrants, and within each stage the three
regions move sequentially. \textcolor{teal}{[Stas] consider the alternative
reduction we discussed to cut computational burden.} Each sub-stage reduces
to a closed-form logit, so no fixed-point iteration is required.
Empirically, we estimate the model in two steps that identify agglomeration
($\gamma$) separately from the action-cost parameters ($\kappa, \phi$).
Nonlinear least squares on the deterministic Cournot quantities pins
$\gamma$. Maximum likelihood with the stage order replayed pins
$(\kappa, \phi)$. For policy, we study innovation and entry subsidies
targeted at a single region. Each is welfare-reducing in isolation, but a
\emph{joint} subsidy tilted toward entry generates a small positive
aggregate effect when agglomeration spillovers are active.
